\documentclass[10pt]{article}

\usepackage[margin=0.72in]{geometry}
\usepackage{amsmath,amssymb}
\usepackage{microtype}

\newcommand{\T}{\mathcal{T}}
\newcommand{\A}{\mathcal{A}}
\newcommand{\Rnn}{\mathbb{R}_{\ge 0}}

\title{\vspace{-1.5em}Catamorphic Weightings for Operational Theories}
\author{}
\date{\vspace{-2em}}

\begin{document}
\maketitle
\vspace{-1.2em}

An \emph{operational theory} consists of a cartesian closed category with finite
limits $\T$, distinguished objects $P$ and $R$ of processes and rewrites,
source and target maps
\[
s,t:R\to P,
\]
and an identity-on-objects functor
\[
F:\A\to\T,
\]
where $\A$ is the free cartesian closed category with finite limits on the
objects of $\T$.  The morphisms added in passing from $\A$ to $\T$ may be
regarded as the constructors and operational rules of a programming language.
The pair $s,t$ makes the global elements of $P$ and $R$ into a quiver, and a
stochastic operational semantics may be obtained by assigning each rewrite a
nonnegative weight, subject to suitable summability conditions on the incoming
and outgoing fibers.

For an ordinary multi-sorted Lawvere theory, such a weighting can be described
catamorphically: replace the rewrite sort by a carrier of weights and interpret
each constructor producing a rewrite by an operation on weights.  The same
idea extends to operational theories, provided that the replacement is
performed through the entire cartesian closed arity structure rather than only
on first-order signatures.  Let $\Omega$ be any set
and choose a finite-limit- and cartesian-closed interpretation
\[
J:\A\to\mathbf{Set}
\]
with $J(R)=\Omega$.  One may, for example, take $J(P)=1$ when process values
themselves are irrelevant to the weighting calculation.  Cartesian closure
then propagates the interpretation through higher-order arities:
\[
J(R\times R)=\Omega\times\Omega,
\qquad
J(R^R)=\Omega^\Omega.
\]
A weighting algebra is a structure-preserving functor
\[
W:\T\to\mathbf{Set}
\]
whose restriction along $F$ is $J$.  For a closed rewrite $r:1\to R$, the map
$W(r):1\to\Omega$ determines its weight $w(r)$.  Thus the desired weighting is
obtained by folding the constructors and operational rules of $\T$ into an
algebra of weight operations.  Higher-order constructors cause no special
difficulty: a constructor $R^R\to R$, for example, is interpreted by a
function $\Omega^\Omega\to\Omega$.  To obtain a genuine catamorphism from a
small collection of operations, however, $\T$ must be given freely, or by
generators and equations, over $\A$; the bare data of an operational theory do
not by themselves provide a recursive decomposition of all morphisms of
$\T$.

In pi-calculus-like theories, the same catamorphism can instead be
\emph{internalized} in the syntax.  Suppose communication is generalized from
name binding to filtered pattern matching, with receivers such as
\[
\texttt{for ([head,\ldots tail] <- x where head < 10) \{ P \}}.
\]
Ordinarily the expression after \texttt{where} is Boolean and determines
whether a successful unification may fire.  Replace Boolean guards by
$\Omega$-valued expressions.  A match producing bindings $b$ then generates
the corresponding communication rewrite with weight $g(b)$, where
$g:B\to\Omega$ is the expression appearing in the \texttt{where} clause.
Boolean filtering embeds as the special case in which false has weight $0$
and true has weight $1$.  Categorically, one may introduce an object $\Omega$
directly into $\T$ together with a morphism
\[
\omega:R\to\Omega,
\]
so that the operational rule generating a communication rewrite also
determines its value under $\omega$.  Additional block syntax can supply
lexically scoped weight functions or policies, allowing a local
\texttt{where} weight to be combined, for example multiplicatively, with an
ambient weight associated with the communication rule.  In this form the
catamorphism is no longer merely an external interpretation of operational
syntax: the language itself contains the algebra from which rewrite weights
are computed.

The remaining summability requirement is global rather than catamorphic.
Local rules may compositionally construct $\omega:R\to\Omega$, while the
conditions
\[
\sum_{s(r)=p} w(r)<\infty
\qquad\text{and}\qquad
\sum_{t(r)=p} w(r)<\infty
\]
must be imposed separately on admissible models, policies, or processes.
Once these hold, normalization of the outgoing weights produces the desired
stochastic transition system.

\end{document}
